\section{The construction: from packaging to a metric}
\label{sec:construction}

This section defines the objects used throughout and proves the basic metric facts.
Nothing here depends on a particular substrate.

\subsection{Micro dynamics, lenses and prototypes}

Let $Z$ be a finite nonempty set of microstates and $P\in\R^{Z\times Z}$ a row-stochastic Markov kernel~\cite{Norris1997}.
Distributions are row vectors, so one step maps $\mu$ to $\mu P$, and the stage parameter $\tau\in\N$ gives the $\tau$-step map $\mu\mapsto\mu P^{\tau}$.

A \emph{lens} is a map $f\colon Z\to X$ onto a finite set $X$ of macro states (nonempty, since $f$ is onto).
Its fibres $f^{-1}(x)$ are the blocks of microstates that the lens does not distinguish; they are the candidate \emph{points} at this resolution.
The lens is encoded by the deterministic coarse matrix $C\in\{0,1\}^{Z\times X}$, $C[z,x]=\one\{f(z)=x\}$, which pushes a micro distribution $\mu$ forward to $\mu C$.

To go back from macro to micro we choose, for each $x\in X$, a \emph{prototype} $u_x$, a probability distribution on $Z$.
The prototypes form the rows of a stochastic matrix $U\in\R^{X\times Z}$.
We always use prototypes supported on their own fibres, $u_x(z)=0$ whenever $f(z)\neq x$.
Then every fibre is nonempty and
\begin{equation}
UC=I_X,
\label{eq:UC}
\end{equation}
so lifting a macro distribution and repackaging it returns the same macro distribution.
Two prototype choices are used: \emph{uniform} on the fibre, and \emph{stationary-conditional}, which restricts a chosen stationary law $\pi$ of $P$ to the fibre and renormalizes; the latter requires $\pi(f^{-1}(x))>0$ for every $x$.

\subsection{Closure and the macro kernel}

Lifting, evolving for $\tau$ steps and repackaging defines two operators:
\begin{equation}
E \coloneqq P^{\tau}CU \in\R^{Z\times Z},
\qquad
K \coloneqq UP^{\tau}C \in\R^{X\times X}.
\label{eq:closure-macro}
\end{equation}
$E$ is the \emph{closure} acting on micro distributions (evolve, repackage, re-lift), and $K$ is the \emph{macro kernel}: $K_{xy}$ is the probability that a walker started from prototype $u_x$ is in block $y$ after $\tau$ steps.
Both are row-stochastic whenever $P$ and $U$ are (\cref{prop:closure-basic}).
This is the simplest form of the operator-rewrite primitive: micro dynamics rewritten in macro variables by inserting a packaging map and a chosen lift.

\subsection{Likelihood costs}
\label{sec:costs}

To obtain an undirected geometry we symmetrize macro transition probabilities into a weight array
\begin{equation}
W\coloneqq \tfrac12\bigl(K+K^{\mathsf T}\bigr),
\end{equation}
which is symmetric with entries in $[0,1]$ (it need not be stochastic).
Fix a probability floor $0<\eta\le 1$ and an edge threshold $\varepsilon_{\mathrm{edge}}\ge 0$.
For $x\neq y$ the edge $\{x,y\}$ is present when $W_{xy}>\varepsilon_{\mathrm{edge}}$, and then its cost is
\begin{equation}
c(x,y)\coloneqq -\log\max\bigl(W_{xy},\eta\bigr)\ \in\ [0,\,-\log\eta].
\label{eq:cost}
\end{equation}
Absent edges have cost $+\infty$.
The floor only caps the cost of very unlikely present edges; it never creates an edge.
Costs are nonnegative because $\max(W_{xy},\eta)\le 1$.

\begin{remark}[Why a floor and not additive smoothing]
The superficially similar rule $-\log(W_{xy}+\eta)$ is negative whenever $W_{xy}+\eta>1$, for instance when $W_{xy}=1$.
A negative undirected edge creates negative-cost cycles, so the infimum over paths can be $-\infty$ and shortest-path algorithms are invalid.
The floor rule~\eqref{eq:cost} avoids this.
In the limit constructions of \cref{sec:constructed} the floor is chosen below every positive edge weight, so it is inactive; a fixed positive floor would eventually clip edges of a refining family whose transition probabilities tend to zero.
\end{remark}

\subsection{Distance as cheapest protocol}

A \emph{protocol} from $x$ to $y$ is a finite path $\gamma=(x=x_0,x_1,\dots,x_k=y)$ in the macro graph, including the empty path when $x=y$.
Its cost is additive, $\mathrm{Cost}(\gamma)=\sum_{i<k}c(x_i,x_{i+1})$, and the induced distance is
\begin{equation}
d(x,y)\coloneqq \inf_{\gamma\colon x\rightsquigarrow y}\mathrm{Cost}(\gamma)\ \in[0,\infty].
\label{eq:path-metric}
\end{equation}
In computation, $d$ is obtained by an all-pairs shortest-path search.
Edges of cost exactly zero must be kept as edges: treating them as absent would disconnect precisely the zero-distance classes that \cref{prop:metric} identifies.

\begin{proposition}[The likelihood path distance]
\label{prop:metric}
For every finite macro kernel $K$, floor $0<\eta\le1$ and threshold $\varepsilon_{\mathrm{edge}}\ge0$:
\begin{enumerate}[label=(\roman*)]
\item $d$ is an extended pseudometric on $X$: $d(x,x)=0$, $d(x,y)=d(y,x)$, and $d(x,z)\le d(x,y)+d(y,z)$, with values in $[0,\infty]$.
\item $d(x,y)<\infty$ for all $x,y$ if and only if the macro graph is connected.
\item Identifying points at distance zero gives an extended metric space whose distances are those of $d$.
It is an ordinary metric space exactly when the graph is connected.
\end{enumerate}
\end{proposition}

\begin{proof}
The empty protocol gives $d(x,x)=0$.
Concatenating a protocol from $x$ to $y$ with one from $y$ to $z$ gives a protocol from $x$ to $z$ whose cost is the sum, which yields the triangle inequality after taking infima; reversing a protocol preserves its cost because $c$ is symmetric.
Nonnegativity of~\eqref{eq:cost} keeps all values in $[0,\infty]$.
Part (ii) is immediate from the definition, since only absent edges have infinite cost.
Part (iii) is the separation quotient of an extended pseudometric.
\end{proof}

\provenance{Lean: weighted protocols and their costs are defined in \lean{PathMetric.lean}, which proves \lean{weightedEdist_self}, \lean{weightedEdist_triangle}, \lean{weightedEdist_symm} and builds \lean{weightedPseudoEMetricSpace} without assuming connectivity, positive costs or separation. \lean{weightedSeparationMetric} applies the separation quotient of \lean{QuotientMetric.lean} to this distance, and \lean{flooredLikelihoodCost_nonneg} proves nonnegativity of~\eqref{eq:cost} for $W\le1$ and $0<\eta\le1$.}

Zero distances between distinct points do occur.
For the two-state deterministic swap with threshold $\varepsilon_{\mathrm{edge}}<1$ (for instance $0$), both off-diagonal probabilities equal one, so both edges are present with cost zero and the two states are at distance zero.
Separation is therefore a property to be checked, not assumed.
The constructions of \cref{sec:constructed} prove explicit positive lower bounds on all edge costs together with connectivity, which gives separation and finiteness directly.

\begin{remark}[Directed costs]
Using $K$ instead of $W$ gives a directed shortest-path distance, an extended quasi-pseudometric: it satisfies the triangle inequality but need not be symmetric, and distinct points may again be at distance zero.
Directed geometries are natural under asymmetric constraints but are not studied here.
\end{remark}
